
\subsection{Natural Language vs. Shell Command Classifier}

\paragraph{Evaluation Approach}{}
The testing for this component focuses on evaluating the accuracy of input
classification across different local \gls{llm} models. A fixed, balanced dataset
is classified into natural language vs.\ shell command inputs.

\paragraph{Testing Collections}{}
For testing the classifier, the manually curated collection\\
\texttt{classifier-100} contains 100 inputs with a balanced mix of natural
language queries and shell command inputs. It is split as follows:
\begin{itemize}
  \item 50 natural language queries (e.g., ``How do I list files in a
    directory?'')
  \item 50 shell command inputs (e.g., ``ls -la /home/user'')
\end{itemize}

\paragraph{Testing Results}{}
Table~\ref{tab:eval-classifier-results} summarizes the classification accuracy
and average per-query latency for each model:

\begin{table}[h]
  \centering
  \scriptsize
  \caption{Binary classification results on \texttt{classifier-100}.}
    \begin{tabular}{lrr}
    \toprule
    \textbf{Model} & \textbf{Accuracy} & \textbf{Avg Latency (s)} \\
    \midrule
    ibm/granite3.1-moe:1b & 53.0\% & 0.278 \\
    ibm/granite3.1-moe:3b & 60.0\% & 0.101 \\
    ibm/granite3.1-dense:2b & 83.0\% & 0.377 \\
    qwen3:0.6b & 79.0\% & 1.183 \\
    qwen3:1.7b & 88.0\% & 2.871 \\
    qwen2.5:0.5b & 50.0\% & 0.180 \\
    qwen2.5:1.5b & 78.0\% & 0.150 \\
    gemma3:1b & 65.0\% & 0.233 \\
    \bottomrule
  \end{tabular}\label{tab:eval-classifier-results}
\end{table}

\paragraph{Key Findings and Limitations}
Evaluation and testing of this component revealed critical limitations that
prevented deployment:

\begin{enumerate}
  \item \textbf{Accuracy Ceiling:} Even the best-performing model
    (qwen3:1.7b) reached only 88\% accuracy on a balanced dataset, while several
    smaller models remained near chance-level performance (50--65\%).
  \item \textbf{Latency vs.\ Accuracy Trade-off:} Higher accuracy required
    larger models with substantially higher latency (up to 2.9s per query),
    which is too slow for interactive shell use.
  \item \textbf{Inconsistent Behavior on Ambiguous Inputs:} Misclassifications
    were frequent on short or context-dependent prompts, which are common in
    CLI workflows.
\end{enumerate}

\paragraph{Conclusion on Classifier}
This component remains a \textbf{proof-of-concept} and was not
pursued further. The fundamental limitation is that small local
\glspl{llm} cannot reliably classify ambiguous inputs in real-time
without either: (a) unacceptable latency through multiple calls, or
(b) extensive hardcoding of program paths that reduce
generalizability, (c) the shell wrapper itself becomes a bottleneck
and maintenance nightmare. The project focus was therefore shifted to
the \gls{rag} component, which proved more practical and effective.

\subsection{Document Retriever}

\paragraph{Evaluation Approach}
The retriever is evaluated on its ability to return relevant documents within
top-k results using Precision@k and MRR, with three automated test variants:
\begin{enumerate}
  \item Document-level retrieval (test1)
  \item Document retrieval with passage-level re-ranking (test2)
      \textit{(retrieving documents first, then re-ranking passages within them, evaluating by passage matches)}
  \item Passage-scope retrieval (test3)
\end{enumerate}
Test1 checks whether the correct document appears in the top-k results. Test2
first retrieves top-k documents and then re-ranks passages within them
(top\_k\_extended=10) to locate the expected section. Test3 indexes and retrieves
passages directly (passage\_size=20 lines), evaluating whether the correct
document is found via passage matches.

\paragraph{Testing Collections}
\begin{itemize}
  \item \texttt{sample-10}: 14 documents, 250 queries
  \item \texttt{sample-100}: 101 documents, 300 queries
\end{itemize}

\paragraph{Testing Results}{}
Tables~\ref{tab:eval-retriever-sample10-doc}--\ref{tab:eval-retriever-sample100-passage}
summarize document-level retrieval (test1), passage-level re-ranking (test2), and
passage-level retrieval (test3). All tests used top\_k=5; test2 uses
top\_k\_extended=10.

\begin{table}[h]
  \centering
  \scriptsize
  \caption{Document-level retrieval on \texttt{sample-10} (test1).}
    \begin{tabularx}{\textwidth}{lYYYYYY}
    \toprule
    \textbf{Embedding Model} & \textbf{Top-1} & \textbf{Top-5} & \textbf{MRR} & \textbf{Avg\newline{}Latency (s)} & \textbf{Index\newline{}Build (s)} & \textbf{Avg Response\newline{}Size (chars)} \\
    \midrule
    ibm/granite-embedding:30m  & 83.6\% & 95.2\%  & 0.881 & 0.071 & 0.348  & 76791 \\
    ibm/granite-embedding:125m & 85.2\% & 92.4\%  & 0.880 & 0.075 & 16.603 & 79930 \\
    ibm/granite-embedding:278m & 84.8\% & 87.6\%  & 0.861 & 0.099 & 2.292  & 77693 \\
    qwen3-embedding:0.6b       & 86.8\% & 100.0\% & 0.916 & 0.074 & 17.503 & 102673 \\
    \bottomrule
  \end{tabularx}\label{tab:eval-retriever-sample10-doc}
\end{table}

\begin{table}[h]
  \centering
  \scriptsize
  \caption{Document retrieval with passage re-ranking on \texttt{sample-10} (test2).}
    \begin{tabularx}{\textwidth}{lYYYYYY}
    \toprule
    \textbf{Embedding Model} & \textbf{Top-1} & \textbf{Top-5} & \textbf{MRR} & \textbf{Avg\newline{}Latency (s)} & \textbf{Index\newline{}Build (s)} & \textbf{Avg Response\newline{}Size (chars)} \\
    \midrule
    ibm/granite-embedding:30m  & 48.4\% & 87.2\% & 0.628 & 0.817 & 1.138  &  5304 \\
    ibm/granite-embedding:125m & 39.2\% & 87.6\% & 0.557 & 1.264 & 1.044  &  5092 \\
    ibm/granite-embedding:278m & 46.4\% & 81.6\% & 0.594 & 1.514 & 1.900  &  5322 \\
    qwen3-embedding:0.6b       & 55.2\% & 94.0\% & 0.685 & 3.777 & 10.322 &  4820 \\
    \bottomrule
  \end{tabularx}\label{tab:eval-retriever-sample10-rerank}
\end{table}

\begin{table}[h]
  \centering
  \scriptsize
  \caption{Passage-level retrieval on \texttt{sample-10} (test3).}
    \begin{tabularx}{\textwidth}{lYYYYYY}
    \toprule
    \textbf{Embedding Model} & \textbf{Top-1} & \textbf{Top-5} & \textbf{MRR} & \textbf{Avg\newline{}Latency (s)} & \textbf{Index\newline{}Build (s)} & \textbf{Avg Response\newline{}Size (chars)} \\
    \midrule
    ibm/granite-embedding:30m &  76.0\% & 98.4\% & 0.863 & 0.045 & 1.7505  & 2589 \\
    ibm/granite-embedding:125m & 84.0\% & 97.2\% & 0.894 & 0.056 & 3.4254  & 2588 \\
    ibm/granite-embedding:278m & 82.0\% & 97.6\% & 0.888 & 0.100 & 5.2866  & 2669 \\
    qwen3-embedding:0.6b &       86.4\% & 96.8\% & 0.903 & 0.081 & 15.6794 & 2493 \\
    \bottomrule
  \end{tabularx}\label{tab:eval-retriever-sample10-passage}
\end{table}

\begin{table}[h]
  \centering
  \scriptsize
  \caption{Document-level retrieval on \texttt{sample-100} (test1).}
    \begin{tabularx}{\textwidth}{lYYYYYY}
    \toprule
    \textbf{Embedding Model} & \textbf{Top-1} & \textbf{Top-5} & \textbf{MRR} & \textbf{Avg\newline{}Latency (s)} & \textbf{Index\newline{}Build (s)} & \textbf{Avg Response\newline{}Size (chars)} \\
    \midrule
    ibm/granite-embedding:30m &  62.7\% & 86.3\% & 0.724 & 0.043 & 2.270   & 70590  \\
    ibm/granite-embedding:125m & 69.7\% & 88.7\% & 0.768 & 0.036 & 1.706   & 62318  \\
    ibm/granite-embedding:278m & 65.7\% & 84.0\% & 0.734 & 0.098 & 1.652   & 81213  \\
    qwen3-embedding:0.6b &       72.7\% & 92.0\% & 0.806 & 0.074 & 128.222 & 127590 \\
    \bottomrule
  \end{tabularx}\label{tab:eval-retriever-sample100-doc}
\end{table}

\begin{table}[h]
  \centering
  \scriptsize
  \caption{Document retrieval with passage re-ranking on \texttt{sample-100} (test2).}
    \begin{tabularx}{\textwidth}{lYYYYYY}
    \toprule
    \textbf{Embedding Model} & \textbf{Top-1} & \textbf{Top-5} & \textbf{MRR} & \textbf{Avg\newline{}Latency (s)} & \textbf{Index\newline{}Build (s)} & \textbf{Avg Response\newline{}Size (chars)} \\
    \midrule
    ibm/granite-embedding:30m &  42.7\% & 72.7\% & 0.530 & 0.595 & 2.132   & 6544 \\
    ibm/granite-embedding:125m & 43.0\% & 77.0\% & 0.539 & 0.802 & 2.451   & 6462 \\
    ibm/granite-embedding:278m & 45.7\% & 72.7\% & 0.549 & 1.192 & 3.219   & 6438 \\
    qwen3-embedding:0.6b &       45.7\% & 79.7\% & 0.571 & 3.745 & 128.563 & 6584 \\
    \bottomrule
  \end{tabularx}\label{tab:eval-retriever-sample100-rerank}
\end{table}

\begin{table}[h]
  \centering
  \scriptsize
  \caption{Passage-level retrieval on \texttt{sample-100} (test3).}
    \begin{tabularx}{\textwidth}{lYYYYYY}
    \toprule
    \textbf{Embedding Model} & \textbf{Top-1} & \textbf{Top-5} & \textbf{MRR} & \textbf{Avg\newline{}Latency (s)} & \textbf{Index\newline{}Build (s)} & \textbf{Avg Response\newline{}Size (chars)} \\
    \midrule
    ibm/granite-embedding:30m &     67.3\% & 91.0\% & 0.771 & 0.045 & 10.879 & 3429 \\
    ibm/granite-embedding:125m &    68.7\% & 89.3\% & 0.765 & 0.040 & 17.602 & 3337 \\
    ibm/granite-embedding:278m &    69.0\% & 88.3\% & 0.762 & 0.109 & 19.135 & 3410 \\
    qwen3-embedding:0.6b &          67.7\% & 87.0\% & 0.746 & 0.089 & 44.709 & 3281 \\
    \bottomrule
  \end{tabularx}\label{tab:eval-retriever-sample100-passage}
\end{table}

\paragraph{Key Findings}
The document retriever component proved robust and practical with all approaches
in terms of accuracy. On \texttt{sample-10}, document-level Top-1 accuracy
ranged from 83.6--86.8\% with Top-5 accuracy up to 100\%. On
\texttt{sample-100}, Top-1 ranged from 62.7\% to 72.7\% with Top-5 up to 92.0\%.
The qwen3-embedding:0.6b model provided the best precision, but its index build
time was substantially higher on the larger corpus (128--134s vs.\ 1.6--2.3s for
Granite embeddings). Though, it is important to note that secound and third
approaches have much smaller output size (avg. 5K characters, which is about 4K
tokens\ vs. avg 80K chars in first approach), which is critical for the
summarization component, as it significantly reduces the latency and increases
the accuracy of the generated answer.

\subsection{End-to-end testing}

\paragraph{Evaluation Approach}
The complete \gls{rag} system (document retrieval + information
summarization) is evaluated end-to-end to assess real-world
performance including document indexing time, query processing, and
user response latency under different hardware configurations and document
corpus sizes. The evaluation focuses on practical usability for \gls{cli}
assistance.

\paragraph{Testing Collection}{}
For end-to-end testing, we use query set (10 queries) and run it against the \texttt{sample-10}, \texttt{sample-100}, and \texttt{sample-1000} corpora.

\paragraph{Testing Results}{}
Table~\ref{tab:eval-21-results}--\ref{tab:eval-32-results} presents sample results from evaluation
for each collection. Values are average per-query latencies in seconds. The
evaluation was conducted on a system with an NVIDIA RTX 4060 8GB GPU. The testing
is conducted for the lightest (qwen2.5:0.5b for generation and
ibm/granite-embedding:30m for retrieval) and the heaviest (qwen3:1.7b for
generation and qwen3-embedding:0.6b for retrieval) models, to provide a range of
expected performance across different hardware configurations.

% \begin{table}[h]
%   \centering
%   \caption{Evaluation results on different corpus sizes text generation on qwen2.5:0.5b}
%   \scriptsize
%   \begin{tabularx}{llll}
%     \toprule
%       \textbf{Name} & \textbf{Approx. input\newline text size (char/tok)} &  \textbf{Avg. Latency (s)} \\
%     \midrule
%     generate-s & <250/180     & None \\
%     generate-m & ~5000/4000   & None \\
%     generate-l & ~60000/50000 & None \\
%     \bottomrule
%   \end{tabularx}\label{tab:eval-11-results}
% \end{table}
%
% \begin{table}[h]
%   \centering
%   \caption{Evaluation results on different corpus sizes text generation on qwen3:1.7b }
%   \scriptsize
%   \begin{tabularx}{lccc}
%     \toprule
%       \textbf{Name} & \textbf{Approx. input\newline text size (char/tok)} & \textbf{Avg. Latency (s)} \\
%     \midrule
%     generate-s & <250/180     & None \\
%     generate-m & ~5000/4000   & None \\
%     generate-l & ~60000/50000 & None \\
%     \bottomrule
%   \end{tabularx}\label{tab:eval-12-results}
% \end{table}

\begin{table}[h]
  \centering
  \caption{Evaluation results on search for different DB approaches on ibm/granite-embedding:30m}
  \scriptsize
    \begin{tabularx}{\textwidth}{lYYY}
    \toprule
    \textbf{DB Collection} & \textbf{Document-level} &
    \textbf{Document-level with passage-level re-ranking} & \textbf{Passage-level} \\
    \midrule
    sample-10   & 1.22 & 1.66 & 1.27 \\
    sample-100  & 1.27 & 1.53 & 1.31 \\
    sample-1000 & 1.37 & 1.64 & 1.27 \\
    \bottomrule
  \end{tabularx}\label{tab:eval-21-results}
\end{table}

\begin{table}[h]
  \centering
  \caption{Evaluation results on search for different DB approaches on qwen3-embedding:0.6b}
  \scriptsize
    \begin{tabularx}{\textwidth}{lYYY}
    \toprule
    \textbf{DB Collection} & \textbf{Document-level} &
    \textbf{Document-level with passage-level re-ranking} & \textbf{Passage-level} \\
    \midrule
    sample-10   & 1.31 & 5.17 & 1.37 \\
    sample-100  & 1.39 & 6.98 & 1.35 \\
    sample-1000 & 0.67 & 16.19 & 0.84 \\
    \bottomrule
  \end{tabularx}\label{tab:eval-22-results}
\end{table}

\begin{table}[h]
  \centering
    \caption{Evaluation results on RAG-search (retrieval + summarization) for different DB approaches on ibm/granite-embedding:30m}
  \scriptsize
    \begin{tabularx}{\textwidth}{lYYY}
    \toprule
    \textbf{DB Collection} & \textbf{Document-level} &
    \textbf{Document-level with passage-level re-ranking} & \textbf{Passage-level} \\
    \midrule
    sample-10   & 3.59 & 2.46 & 1.82 \\
    sample-100  & 3.18 & 1.94 & 1.88 \\
    sample-1000 & 3.36 & 2.09 & 1.94 \\
    \bottomrule
  \end{tabularx}\label{tab:eval-31-results}
\end{table}

\begin{table}[h]
  \centering
    \caption{Evaluation results on RAG-search (retrieval + summarization) for different DB approaches on qwen3-embedding:0.6b and qwen3:1.7b}
  \scriptsize
    \begin{tabularx}{\textwidth}{lYYY}
    \toprule
    \textbf{DB Collection} & \textbf{Document-level} &
    \textbf{Document-level with passage-level re-ranking} & \textbf{Passage-level} \\
    \midrule
    sample-10   & 20.85 & 6.68 & 5.63 \\
    sample-100  & 34.24 & 8.62 & 5.89 \\
    sample-1000 & 26.21 & 10.46 & 6.27 \\
    \bottomrule
  \end{tabularx}\label{tab:eval-32-results}
\end{table}

\begin{table}[h]
  \centering
    \caption{Evaluation results on database build time for different DB approaches on ibm/granite-embedding:30m}
  \scriptsize
    \begin{tabularx}{\textwidth}{lYYY}
    \toprule
    \textbf{DB Collection} & \textbf{Document-level} &
    \textbf{Document-level with passage-level re-ranking} & \textbf{Passage-level} \\
    \midrule
    sample-10   & 0.71 & 0.78 & 1.63 \\
    sample-100  & 1.58 & 1.59 & 7.18 \\
    sample-1000 & 13.10 & 12.74 & 109.03 \\
    \bottomrule
  \end{tabularx}\label{tab:eval-41-results}
\end{table}

\begin{table}[h]
  \centering
    \caption{Evaluation results on database build rime for different DB approaches on qwen3-embedding:0.6b}
  \scriptsize
    \begin{tabularx}{\textwidth}{lYYY}
    \toprule
    \textbf{DB Collection} & \textbf{Document-level} &
    \textbf{Document-level with passage-level re-ranking} & \textbf{Passage-level} \\
    \midrule
    sample-10   & 9.94 & 10.12 & 5.95 \\
    sample-100  & 145.65 & 174.72 & 40.59 \\
    sample-1000 & 1606.91 & 1604.53 & 643.04 \\
    \bottomrule
  \end{tabularx}\label{tab:eval-42-results}
\end{table}

\paragraph{Key Findings}
Document-level retrieval and document-level retrieval with passage reranking
show the best performance when using small embedding models like
ibm/granite-embedding:30m. The latency of runtime passage-level reranking is low
enough to perform better, than passage-level approach, but keeping the output
size small and accuracy the same. In contrast, when using bigger embedding
models, like qwen3-embedding:0.6b, the passage-level approach outperforms the
other approaches, as the computation time needed for vector creation for
the whole document or runtime reranking becomes too expensive.

\subsection{Regression tests}
The software regression tests that were implemented to
ensure the stability and reliability of the ShAI system during development and
future maintenance. The tests cover key paths of the application, including
document retrieval, summarization, and end-to-end query processing.
Along with the software regression tests, we also implemented a test regression
suite, which ensures that with any change of the system, none of the tests are
broken (especially statistical tests, which usually need a lot of time to run).
This allows us to maintain the stability of the system with ease, and to quickly
identify any issues that may arise during development.
